\section{The Execution Engine}
\label{sec:engine}

\sys turns an analysis into a fold over compressed traversal columns. This
section describes the substrate every analysis shares: the decode primitive
(\S\ref{subsec:eng-stream}), what stays resident while a query runs
(\S\ref{subsec:eng-memory}), how the work is parallelized
(\S\ref{subsec:eng-parallel}), and the interface an analysis is written against
(\S\ref{subsec:eng-api}). Section~\ref{sec:deconstruct} then develops the hardest
analysis in full, and \S\ref{subsec:eng-instances} sketches the rest.

\subsection{The Streaming Primitive}
\label{subsec:eng-stream}

Everything reduces to one operation:
\[
  \textsc{StreamNodes}(\textit{slice},\ \textsc{Visit}),
\]
which emits the orientation-signed node IDs of one traversal, in order, to a
caller-supplied visitor. A \textit{slice} is a zero-copy view into the
Zstd-decompressed traversal column, described by a pointer, an encoded length, and
the original expanded length: no copy is made and no traversal is materialized to
create one.

Algorithm~\ref{alg:stream} gives the decoder. For each encoded symbol it performs
the range check of Section~\ref{sec:container}. A symbol outside the rule range is
a literal and is emitted. A symbol inside it indexes the rule arrays at
$|e| - \textit{min\_rule}$ and expands its two children, recursively, carrying a
\emph{reverse} flag that is folded with each child's sign, so that a negative rule
reference visits its children in reverse order with orientations flipped. The
inverse delta transform is fused into the emit step rather than run as a separate
pass: for the common single-round case the decoder carries a running prefix sum
and calls the visitor immediately, so a node ID is reconstructed and consumed
without ever being written to memory.

\begin{algorithm}[t]
\small
\caption{Streaming traversal expansion}
\label{alg:stream}
\begin{algorithmic}[1]
  \REQUIRE slice $E$; rule range $[m,M)$; rule arrays $R_1,R_2$;
           delta round $D\in\{0,1\}$; visitor \textsc{Visit}
  \STATE $\textit{prev}\gets 0$
  \STATE \textbf{define} \textsc{Emit}($v$):
         \textbf{if} $D{=}1$ \textbf{then} $\textit{prev}\mathrel{+}=v$;
         \textsc{Visit}(\textit{prev}) \textbf{else} \textsc{Visit}($v$)
  \STATE \textbf{define} \textsc{Expand}($r$, \textit{rev}):
  \STATE \quad $(a,b)\gets(R_1[r{-}m],R_2[r{-}m])$
  \STATE \quad $(c_1,c_2)\gets \textit{rev}\ ?\ (b,a):(a,b)$
         \COMMENT{reverse swaps child order}
  \STATE \quad \textbf{for} $c \in (c_1,c_2)$:
  \STATE \qquad \textbf{if} $|c|\in[m,M)$ \textbf{then}
         \textsc{Expand}($|c|$, $\textit{rev}\oplus(c{<}0)$)
  \STATE \qquad \textbf{else} \textsc{Emit}($\textit{rev}\ ?\ {-}c:c$)
  \STATE \textbf{for} $e \in E$ in order:
  \STATE \quad \textbf{if} $|e|\in[m,M)$ \textbf{then} \textsc{Expand}($|e|$, $e{<}0$)
         \textbf{else} \textsc{Emit}($e$)
\end{algorithmic}
\end{algorithm}

Three properties of this loop are what make the rest of the system work. Its cost
is linear in the number of emitted elements, with no dependence on file size or on
how many other traversals exist. Rule resolution is an array index, not a hash
probe, and recursion depth is bounded by the number of grammar rounds, which is a
small constant. And the visitor is a template parameter, so the per-node fold is
inlined into the decode loop and the compiler sees a single fused loop rather than
a decoder handing elements to an opaque callback. In effect the query's operator
is compiled into the decompressor.

Higher delta rounds are rare and are handled by a fallback that materializes a
single traversal, applies the required inverse passes, and visits the result. This
preserves correctness for every container while keeping the fast path free of
branches for cases that do not occur in practice.

\subsection{Memory Model}
\label{subsec:eng-memory}

The claim that distinguishes this design from a fast decompressor is a claim about
residency. While a query runs, at most four things are in memory:

\begin{enumerate}[noitemsep,topsep=3pt,leftmargin=*]
  \item \textbf{The compressed traversal store.} Zstd-decompressed but
        \emph{still grammar-compressed} P/W symbol streams, their length columns,
        and the rule arrays. This is the key point: entropy decoding is undone,
        the grammar is not.
  \item \textbf{An optional rule-leaf cache}, bounded by an explicit byte budget.
  \item \textbf{The query's accumulator}, whose shape depends on the analysis.
  \item \textbf{The segment DNA payload}, only for analyses that must spell
        sequence.
\end{enumerate}

Memory therefore scales with the compressed columns and with the accumulator, and
not with the expanded step list, which is the object every baseline builds. For
the analyses whose accumulator is $O(\text{nodes})$ or $O(\text{windows})$, both
terms are small fractions of the input and peak RSS lands well below the
uncompressed GFA. For the analyses whose accumulator is a join index, it does not,
and Section~\ref{sec:discussion} treats that case honestly.

\paragraph{Rule-leaf cache.}
Lazy expansion avoids materialization but can repeat work, since a rule that
appears in many haplotypes is re-expanded at each occurrence. \sys therefore
offers a bounded \texttt{RuleLeafCache} that stores the fully expanded forward
leaf sequence of any rule whose expansion fits a user-specified budget. Reverse
expansion is derived by scanning the cached leaves backward and negating signs, so
no reverse copy is stored (P5 again). Construction is bottom-up over rule IDs: a
rule is admitted only if its children are already expanded or representable
lazily, and only if its projected leaf sequence fits the remaining budget; rules
that do not fit fall back to recursion, so correctness holds at every budget
including zero. The cache is built once and shared read-only across threads.

The budget is the one explicit space-time dial in the engine, and it is the
mechanism by which a user chooses where to sit on the materialization spectrum.
Zero recovers fully lazy expansion; an unbounded budget approaches full
materialization. Only the multi-pass analyses default to a non-trivial budget,
since analyses that stream each traversal once have nothing to reuse. We measure
the dial's sensitivity in Section~\ref{subsec:eval-cache}.

\subsection{Parallel Execution}
\label{subsec:eng-parallel}

Record independence (P4) makes the decode path embarrassingly parallel. The
traversal arrays, rulebook, DNA payload, and rule-leaf cache are all read-only
after setup, so threads decoding distinct traversals never synchronize. Because
expanded traversal lengths vary by orders of magnitude across haplotypes, the
parallel loops use dynamic scheduling.

All synchronization is thus pushed into the fold, never the decode, and each
analysis picks one of three merge strategies according to its accumulator: a
shared accumulator updated with fine-grained atomics, private per-thread state
concatenated after the parallel region, or a parallel join over a node-major
index. Where a merge strategy could make output depend on thread interleaving, the
analysis restores determinism explicitly; \cmd{deconstruct}, for instance, sorts
each site's observations by source slice so its VCF is byte-identical at any
thread count.

\subsection{Writing an Analysis}
\label{subsec:eng-api}

\sys is meant to be extended, so the substrate is exposed as a small library
rather than baked into each command. It provides container loading that
decompresses the requested blocks and exposes traversals as zero-copy slices; the
templated \textsc{StreamNodes} decoder; the optional rule-leaf cache; and a single
grouping facility that maps each traversal to a key at one of four granularities
(per record, per sample, per sample-haplotype, or per sample-haplotype-sequence)
with shared key functions for P and W lines, so that every analysis aggregates
haplotypes identically.

An analysis is then a fold following one skeleton: open the container, assign each
traversal a group key, stream slices into an analysis-specific visitor under a
parallel loop, merge per-thread accumulators, emit. Every command in
Table~\ref{tab:commands} is written this way and differs from the others only in
its accumulator and its output, not in how traversals are decoded. In our
experience adding an analysis is a few hundred lines, most of which is output
formatting.

\subsection{The Analyses}
\label{subsec:eng-instances}

\sys implements six analyses, each reproducing the output of an established tool.
They fall into two regimes, which differ in what the accumulator costs.

\emph{Path-iterative} folds stream each traversal once into an accumulator that
does not grow with the number of haplotypes. \cmd{growth} computes a
Panacus-style union growth curve: threads reduce into a shared coverage vector
under a per-thread last-seen stamp that counts each node at most once per group,
and the vector is folded into a histogram from which the expected number of nodes
observed in a random cohort of size $k$ follows in closed form, evaluated as a
stable running product rather than by forming binomial coefficients that would
overflow at pangenome scale. \cmd{depth} reproduces \cmd{odgi depth} in one pass
with per-node atomic counters. \cmd{stats} is the degenerate case: it answers
from the metadata columns alone and decodes no traversals at all, which is
possible only because of column independence (P3).

\emph{Node-join} workloads instead record which groups cover each node and join
along shared nodes. \cmd{pav} computes a window-by-group presence/absence matrix
over BED intervals on reference paths, in three passes: decode every traversal
into a per-slice node set while capturing reference node order as a side effect;
union those sets per group and flatten them into a node-to-groups index in
compressed-sparse-row form, chosen over a per-node container to eliminate
per-node header overhead at hundreds of millions of nodes; then sweep each
reference contig, mapping nodes to genomic intervals through the segment length
column and intersecting against pre-sorted BED ranges. \cmd{similarity}
generalizes the same join from reference windows to all group pairs, accumulating
an upper-triangular intersection matrix, and computes in double precision to avoid
the integer overflow the reference tool incurs at genome-scale group lengths.

The node-join regime is where the memory story changes, because the CSR index and
output matrix are ordinary uncompressed structures whose size tracks the graph
rather than the query. That is the origin of the limitation in
Section~\ref{sec:discussion}.

\begin{table}[t]
\centering
\footnotesize
\setlength{\tabcolsep}{3pt}
\caption{The six analyses, all built on one decode substrate. ``Resident state''
excludes the compressed traversal store common to every analysis. Node counts
dominate group and window counts at pangenome scale.}
\label{tab:commands}
\begin{tabular}{@{}llccl@{}}
\toprule
Analysis & Reference tool & Pass & Resident state & Regime \\
\midrule
\cmd{deconstruct} & \cmd{vg deconstr.} & 1{+}1 & $O(\text{sites}){+}$DNA & path-iter. \\
\cmd{growth}      & Panacus              & 1 & $O(\text{nodes})$ & path-iter. \\
\cmd{depth}       & \cmd{odgi depth}     & 1 & $O(\text{nodes})$ & path-iter. \\
\cmd{stats}       & \cmd{odgi stats}     & 0 & $O(1)$ & path-iter. \\
\cmd{pav}         & \cmd{odgi pav}       & 3 & $O(n{+}w{\times}g)$ & node-join \\
\cmd{similarity}  & \cmd{odgi similar.} & 2{+}1 & $O(n{+}g^2)$ & node-join \\
\bottomrule
\end{tabular}
\end{table}
